\section{REST-API-Referenz}
\label{sec:api}

SignalReport stellt eine REST-API auf Port 4567 bereit, die sowohl vom integrierten
Web-Interface als auch von externen Tools genutzt werden kann. Alle Antworten verwenden
das JSON-Format.

\subsection{Authentifizierung}

Wenn die Authentifizierung aktiviert ist, erfolgt der Login über ein
Challenge-Response-Verfahren. Alle nachfolgenden API-Aufrufe werden per
Session-Cookie (\texttt{SR\_SESSION}) authentifiziert.

\begin{lstlisting}[language=bash, caption={Beispiel: Login via Challenge-Response}]
# 1. Nonce holen
curl http://localhost:4567/api/auth/nonce
# -> {"nonce":"a1b2c3..."}

# 2. Login (challengeResponse = SHA-256(SHA-256(passwort) + nonce))
curl -X POST http://localhost:4567/api/auth/login \
  -H "Content-Type: application/json" \
  -d '{"username":"admin","nonce":"a1b2c3...","challengeResponse":"d4e5f6..."}'
# -> {"token":"abc123...","role":"admin"}

# 3. Folgende Requests mit Cookie
curl -b "SR_SESSION=abc123..." http://localhost:4567/api/measurements
\end{lstlisting}

Schreibende Endpunkte (POST) erfordern Admin-Rechte (Admin-Session).

\subsection{Messdaten}

\begin{table}[h]
\centering
\small
\begin{tabular}{|l|l|p{7cm}|}
\hline
\textbf{Methode} & \textbf{Pfad} & \textbf{Beschreibung} \\
\hline
GET & \texttt{/api/measurements} & Letzte Messungen abrufen.
  Parameter: \texttt{limit} (Anzahl, Standard: 50) \\
\hline
GET & \texttt{/api/statistics} & Statistiken für einen Messtyp.
  Parameter: \texttt{type} (PING|DNS|HTTP), \texttt{hours} (Zeitraum) \\
\hline
GET & \texttt{/api/hourly} & Stündliche Durchschnittswerte für Heatmap.
  Parameter: \texttt{type}, \texttt{days} \\
\hline
\end{tabular}
\caption{API-Endpunkte: Messdaten}
\end{table}

\textbf{Beispiel-Antwort} (\texttt{/api/measurements?limit=1}):
\begin{lstlisting}[language=json]
[{
  "timestamp": "2026-03-15T14:32:05Z",
  "target": "8.8.8.8",
  "latencyMs": 23.45,
  "success": true,
  "type": "PING",
  "localIPv4": "192.168.1.100",
  "externalIPv4": "85.182.xxx.xxx",
  "hostHash": "a3f8b2c1..."
}]
\end{lstlisting}

\subsection{Berichte und Export}

\begin{table}[h]
\centering
\small
\begin{tabular}{|l|l|p{7cm}|}
\hline
\textbf{Methode} & \textbf{Pfad} & \textbf{Beschreibung} \\
\hline
GET & \texttt{/api/report} & PDF-Bericht generieren.
  Parameter: \texttt{hours} (24, 168, 8760).
  Antwort: PDF-Datei (application/pdf) \\
\hline
GET & \texttt{/api/export/csv} & CSV-Export.
  Parameter: \texttt{hours} (Zeitfilter), \texttt{all=true} (alle Daten),
  \texttt{type} (Messtyp-Filter).
  Antwort: CSV-Datei (text/csv, Semikolon-getrennt) \\
\hline
\end{tabular}
\caption{API-Endpunkte: Berichte und Export}
\end{table}

\subsection{DNS-Benchmark}

\begin{table}[h]
\centering
\small
\begin{tabular}{|l|l|p{7cm}|}
\hline
\textbf{Methode} & \textbf{Pfad} & \textbf{Beschreibung} \\
\hline
GET & \texttt{/api/dns/servers} & Liste aller konfigurierten DNS-Server abrufen \\
\hline
POST & \texttt{/api/dns/benchmark} & DNS-Benchmark starten.
  Parameter: \texttt{hostname}.
  Fragt alle DNS-Server parallel ab (Virtual Threads).
  30-Sekunden-Cooldown. \\
\hline
GET & \texttt{/api/dns/statistics} & DNS-Statistik (Anzahl Server, Regionen, Provider) \\
\hline
\end{tabular}
\caption{API-Endpunkte: DNS-Benchmark}
\end{table}

\subsection{IP-Tracking und Host-Info}

\begin{table}[h]
\centering
\small
\begin{tabular}{|l|l|p{7cm}|}
\hline
\textbf{Methode} & \textbf{Pfad} & \textbf{Beschreibung} \\
\hline
GET & \texttt{/api/ip-changes} & Letzte IP-Änderungen.
  Parameter: \texttt{limit} (Standard: 50) \\
\hline
GET & \texttt{/api/ip-statistics} & Statistik über IP-Änderungen
  pro Host (Anzahl, erster/letzter Wechsel) \\
\hline
GET & \texttt{/api/hosts} & Liste aller registrierten Hosts \\
\hline
GET & \texttt{/api/host/current} & Aktueller Host: Hostname, Hash, IPv4/IPv6 \\
\hline
\end{tabular}
\caption{API-Endpunkte: IP-Tracking}
\end{table}

\subsection{Konfiguration}

\begin{table}[h]
\centering
\small
\begin{tabular}{|l|l|p{7cm}|}
\hline
\textbf{Methode} & \textbf{Pfad} & \textbf{Beschreibung} \\
\hline
GET & \texttt{/api/config/current} & Aktuelle Konfiguration abrufen
  (Messziele, Intervall, Maintenance, Benutzer-Info) \\
\hline
POST & \texttt{/api/config/update} & Konfiguration aktualisieren.
  Body: JSON mit ping, dns, http, intervalSeconds, maintenance, userInfo \\
\hline
GET & \texttt{/api/theme/settings} & Theme-Einstellungen abrufen \\
\hline
POST & \texttt{/api/theme/settings} & Theme ändern.
  Body: \texttt{\{``darkMode'': true\}} \\
\hline
\end{tabular}
\caption{API-Endpunkte: Konfiguration}
\end{table}

\subsection{Authentifizierung und Setup}

\begin{table}[h]
\centering
\small
\begin{tabular}{|l|l|p{7cm}|}
\hline
\textbf{Methode} & \textbf{Pfad} & \textbf{Beschreibung} \\
\hline
GET & \texttt{/api/auth/nonce} & Einmal-Nonce für Challenge-Response
  generieren (60\,s gültig) \\
\hline
POST & \texttt{/api/auth/login} & Login via Challenge-Response.
  Body: \texttt{\{username, nonce, challengeResponse\}} \\
\hline
POST & \texttt{/api/auth/logout} & Session beenden (Cookie wird gelöscht) \\
\hline
GET & \texttt{/api/auth/status} & Auth-Status abfragen (aktiviert/deaktiviert) \\
\hline
POST & \texttt{/api/setup/complete} & Ersteinrichtung abschließen.
  Body: \texttt{\{adminPasswordHash, enableAuth, userPasswordHash\}} \\
\hline
POST & \texttt{/api/auth/enable} & Auth aktivieren (Challenge-Response + User-Passwort-Hash) \\
\hline
POST & \texttt{/api/auth/disable} & Auth deaktivieren (Challenge-Response erforderlich) \\
\hline
POST & \texttt{/api/auth/change-admin} & Admin-Passwort ändern
  (Challenge-Response für altes PW + neuer Hash) \\
\hline
\end{tabular}
\caption{API-Endpunkte: Authentifizierung}
\end{table}

\subsection{Push-Benachrichtigungen}

\begin{table}[h]
\centering
\small
\begin{tabular}{|l|l|p{7cm}|}
\hline
\textbf{Methode} & \textbf{Pfad} & \textbf{Beschreibung} \\
\hline
GET & \texttt{/api/push/settings} & Push-Einstellungen abrufen \\
\hline
POST & \texttt{/api/push/settings} & Push-Einstellungen ändern
  (Schwellwert, Anzahl) \\
\hline
\end{tabular}
\caption{API-Endpunkte: Push-Benachrichtigungen}
\end{table}

\subsection{HTTP-Statuscodes}

\begin{table}[h]
\centering
\begin{tabular}{|l|p{9cm}|}
\hline
\textbf{Code} & \textbf{Bedeutung} \\
\hline
200 & Erfolgreiche Anfrage \\
302 & Weiterleitung (z.\,B.\ zum Setup) \\
401 & Authentifizierung erforderlich oder fehlgeschlagen \\
403 & Keine Berechtigung (z.\,B.\ User versucht Admin-Aktion) \\
500 & Interner Serverfehler (Details im Response-Body) \\
\hline
\end{tabular}
\caption{Verwendete HTTP-Statuscodes}
\end{table}
